This Rapid Multivocal Review highlights a rapidly expanding landscape of Generative AI tools applicable to software construction, as detailed throughout Section \ref{sec:results} where we systematically categorize and analyze these tools. Section \ref{sec:futurework} critically addresses inherent validity threats, including scope definition, objectivity in data extraction, and researcher bias. Despite significant advancements, overall, we emphasize ongoing challenges in integrating these tools into reliable software development workflows, as some of these tools are still very bare bones and often require extensive environment configuration, which is probably why ChatGPT ends up being the tool that appears the most, due to being easier for the end user. Section \ref{sec:futurework} also outlines future research directions, advocating for the extension of repository searches, the employment of advanced filtering techniques, and the continuous validation of the applicability and maturity of tools. Ultimately, this work provides a structured, evidence-based foundation that encourages deeper academic-industry collaborations to exploit the transformative potential of Generative AI fully.

% \begin{table}[h!]
%     \centering
%     \caption{Examples of Tools by Software Engineering Activity}
%     \begin{tabular}{|l|l|}
%         \hline
%         \textbf{Activity} & \textbf{Example Tools} \\
%         \hline
%         Code Generation & Copilot, Codeium, Ghostwriter \\
%         Test Creation & TestPilot, GPT-TestGen \\
%         Refactoring & StarCoder, CodeT5+ \\
%         Requirement Parsing & Software Engineer GPT, NL2Code \\
%         Bug Fixing & Calico, ASPIRE \\
%         \hline
%     \end{tabular}
%     \label{tab:tools_by_activity}
% \end{table}